% template taken from https://www.dickimaw-books.com/latex/admin/html/leafletcls.shtml

\documentclass[a4paper,12pt,notumble]{leaflet}

\usepackage[utf8]{inputenc}
\usepackage[T1]{fontenc}
\usepackage{tgadventor}

\usepackage{wasysym}
\usepackage{xcolor}
\usepackage{graphicx}
\usepackage{tikz}
\usepackage[utf8]{inputenc} 
\usetikzlibrary{calc}

\renewcommand*{\sectfont}{\scshape}
\renewcommand*{\descfont}{\bfseries\scshape}

\definecolor{rwth-blue}{HTML}{00549F}
\definecolor{rwth-lblue}{HTML}{8EBAE5}
\definecolor{rwth-llblue}{HTML}{C7DDF2}

\AddToBackground{1}{%
	\put (0,0) {
		\begin{tikzpicture}
			%\path[fill=Thistle1,draw=Thistle4,double=Orchid1,line width=2pt] (0,0) rectangle ($(\paperwidth, \paperheight)-(4pt,4pt)$);
			\shadedraw[inner color=white,outer color=gray!20!white,draw=rwth-blue] (0,0) rectangle ($(\paperwidth, \paperheight)-(4pt,0pt)$);
		\end{tikzpicture}
	}
}

\AddToBackground{3}{%  Background of a small page
	\put(0,0){
		\textcolor{rwth-llblue}{\rule{\paperwidth}{\paperheight}}
	}
}
\AddToBackground{4}{%  Background of a small page
	\put(0,0){
		\textcolor{rwth-llblue}{\rule{\paperwidth}{\paperheight}}
	}
}

%\CutLine{6}

%%% Document start -------------------------------------------------------------------------
\begin{document}

%%%%%%%%%%%%%%%%%%
%%% Title page %%%
%%%%%%%%%%%%%%%%%%
\begin{center}

\includegraphics[height=2cm]{../Logos/Cosima18.png}

\bfseries\Huge
Der Gestikulaser

\vfill

\normalsize

\begin{tabular}{c}
	Ein Projekt im Rahmen des \\
	COSIMA-Wettbewerbs 2018
\end{tabular}

\vfill
\includegraphics[scale=0.25]{../Logos/gestikulaser.png}
\vfill

\begin{tabular}{c}
	entwickelt von \\ \\
	Christoph Behr \\
	Cailing Fu \\
	Nicole Grubert \\
	Daniel Wolff
\end{tabular}

\vfill

\includegraphics[height=2cm]{../Logos/TOS.png}
%\includegraphics[height= 1.5cm]{../Logos/ILT.png}

\end{center}

\newpage
\raggedright


%%%%%%%%%%%%%%%%%%%
%%% Second page %%%
%%%%%%%%%%%%%%%%%%%

\noindent
\begin{minipage}[c][0.45\textheight][t]{\textwidth}
	\section{Der Gestikulaser}

	Der Gestikulaser ist ein intelligentes System zur Gestenerkennung. Sie basiert auf Lichtreflextionen der Hand, wodurch benutzerdefinierte Gesten gelernt und
	erkannt werden sollen. 

	\begin{flushright}
		\includegraphics[scale=0.2]{../Logos/gestikulaser.png}
	\end{flushright}
\end{minipage}

\textcolor{rwth-lblue}{\noindent\rule{\textwidth}{4pt}}

\noindent
\begin{minipage}[c][0.45\textheight][t]{\textwidth}
	\begin{itemize}
		\item Intelligenzes System zur Gestensteuerung
		\item Tag und Nacht verwendbar
		\item Individuell auf die Geste des Nutzers abgestimmt
		\item Portabel durch innovatives Stecksystem
		\item Schnittstelle für jedes Bluetooth Endgerät
	\end{itemize}
\end{minipage}

\newpage


%%%%%%%%%%%%%%%%%%
%%% Third page %%%
%%%%%%%%%%%%%%%%%%

% Macht eine Doppelseite
\makebox[\linewidth][l]{\begin{minipage}[c][\textheight][s]{2.2\linewidth}

\section{Funktionsweise} 
\textcolor{red}{Wie benutzt man den Gestikulaser? Was wird derzeit damit gesteuert?}

\end{minipage}}

% Eine leere Seite (nötig wegen der Doppelseite)
\newpage
\thispagestyle{empty}
\quad 
\newpage


%%%%%%%%%%%%%%%%%%
%%% Fifth page %%%
%%%%%%%%%%%%%%%%%%

\noindent
\begin{minipage}[c][0.45\textheight][t]{\textwidth}
	\section{Unsere Zukunftsvision}
	\textcolor{red}{Was kann noch verbessert werden und was ist unsere Vorstellung?}
\end{minipage}

\textcolor{rwth-lblue}{\noindent\rule{\textwidth}{4pt}}

\noindent
\begin{minipage}[c][0.45\textheight][t]{\textwidth}
	\begin{itemize}
		\item Automatisches Lernen dank eines Sensorhandschuhs
		\item Auf beliebige Größe erweiterbar
		\item Individuelles Muster dank beliebig kombinierbares Stecksystem
	\end{itemize}
\end{minipage}

\newpage


%%%%%%%%%%%%%%%%%%
%%% Sixth page %%%
%%%%%%%%%%%%%%%%%%

\noindent
\begin{minipage}[c][0.45\textheight][t]{\textwidth}
	\section{Das Team}
	\begin{description}
		\item[Projektleitung] Cailing Fu
		\item[Teammitglieder] Christoph Behr, Nicole Grubert, Daniel Wolff
	\end{description}

	\section{Kontakt}
	\begin{description}
		\item[Kontaktperson] Cailing Fu 
		\item[Adresse] Fraunhofer-Institut für Lasertechnik, Steinbachstraße 15, 52074 Aachen
		\item[Telefon] +49 176 43404702
		\item[E-Mail] cailing.fu@rwth-aachen.de
	\end{description}
\end{minipage}

\textcolor{rwth-lblue}{\noindent\rule{\textwidth}{4pt}}

\noindent
\begin{minipage}[c][0.45\textheight][t]{\textwidth}
	\begin{center}
		\section{Sponsoren und Unterstützer}
		\vspace{1cm}
		\includegraphics[height= 1cm]{../Logos/ILT.png} \\ \vspace{0.3cm}
		\includegraphics[height= 1.5cm]{../Logos/Wuerth.png} \\ \vspace{0.3cm}
		\includegraphics[height= 1cm]{../Logos/aconity.png} \\ \vspace{0.3cm}
		\includegraphics[height= 1.5cm]{../Logos/BMBF.jpg} \\ \vspace{0.3cm}
		%\includegraphics[height= 0.5cm]{../Logos/Faulhaber.png} \\ \vspace{0.3cm}
		%\includegraphics[height= 1.5cm]{../Logos/micronit.jpg} \\ \vspace{0.3cm}
		\includegraphics[height= 1cm]{../Logos/electronica.png} \\ \vspace{0.3cm}
		\includegraphics[height= 0.8cm]{../Logos/VDE.jpg}
	\end{center}
\end{minipage}

\end{document}